\section{Introduction}
\label{sec:intro}
The International Roughness Index (IRI) is the standard summary of longitudinal road roughness~\cite{sayers1986tp45,astm_e1926}.
It is computed from a longitudinal profile measured by a laser profilometer, an instrument too costly to survey minor roads often.
Production cars that already carry accelerometers could monitor roughness continuously~\cite{sattar2018review,levenberg2021lira}, but their vibration depends on the vehicle, the sensor mounting and the speed.
The mapping from vibration to IRI must therefore be established for each car type from paired data, drives of that car over road sections whose profile has been measured~\cite{sayers1986tp46}, and collecting such drives is the expensive step.

Vehicle simulation offers another source of pairs: driving a vehicle model over profiles that road agencies already survey yields vibration with exact IRI labels, at any speed and for many vehicle parameters.
Convolutional networks trained on simulated responses of four vehicle types estimate IRI in simulation~\cite{jeong2020cnn}; regressors trained on simulated responses were tested with a probe vehicle for the weighted longitudinal profile~\cite{nitsche2014wlp}, as described in~\cite{bajic2021ml}; networks trained on a full-vehicle model over slightly modified measured profiles were applied to that vehicle's measured responses~\cite{ngwangwa2014exp}; and simulation-trained classifiers of roughness grades were tested on a real vehicle~\cite{speedadaptive2025}.
Calibrated vehicle models instead invert car vibration into road profiles~\cite{skar2024inversion,agebjar2025fusion} or derive an IRI--vibration relation by simulation~\cite{wu2023calib}, and learned estimators trained on real drives~\cite{sang2024iri,jeongjo2024irinet}, including on the data used here~\cite{bajic2021ml}, need paired data for every car type.
In robotics, randomised simulation~\cite{tobin2017dr,peng2018dr} adapted with a few real roll-outs~\cite{chebotar2019simopt} is the standard route from simulation to reality.
What has not been shown, for continuous IRI from a production car on held-out roads, is which parts of a simulation-trained estimator transfer, which simulator choices set its accuracy, and how much paired real data must accompany the simulation.

We answer these questions on LiRA-CD~\cite{skar2023lira}, public data of one car model on four routes with profilometer references.
Estimators trained on real drives are compared with the same estimators trained only on simulated drives over the measured profiles of the other routes, under nine sensor and vehicle models of the simulator, and with the simulation combined with $k$\,km of paired drives (Fig.~\ref{fig:pipeline}).
Our contributions are:
\begin{itemize}
\item We show that ranking segments within a route is not the bottleneck: an untrained, speed-normalised vibration RMS ranks them within \rpAbsMax{} in Spearman correlation of the best trained estimators, real or simulated (Section~\ref{sec:res_zero}).
\item We show that the absolute level is the bottleneck, and that it depends on how well the simulated sensor path reproduces the car's vibration above 10\,Hz and on the simulated vehicle population (Section~\ref{sec:res_mech}).
\item We quantify the calibration requirement: about 5\,km of paired drives, with one pass per segment, fix the level, and joint training on simulated and real drives stays better than real-only training up to 20\,km for every simulator variant (Section~\ref{sec:res_value}).
\item We show that simulation must be driven by measured profiles: game-engine road meshes yield an implausible IRI ground truth and do not transfer (Section~\ref{sec:gameengine}).
\end{itemize}
Code, configurations and predictions will be released with the published version.
